\subsection{Equilibria and the three returning strips}

To continue the comparison paths between an initial input and an outer
cycle, we must exclude an equilibrium that could trap an intervening
orbit. We do this first, then fill the upper and two lower strips by
trajectories.

\begin{lemma}[Equilibria enclosed by counted cycles]
\label{rw:fibre:enclosed-equilibria}
On a common small original parameter neighborhood, the Jordan disk
of every counted cycle contains the origin and no other equilibrium.
The section points of any two counted cycles belong to one interval
of the actual return domain.
\end{lemma}
\begin{proof}
The origin has index one and remains the only equilibrium in a
fixed inner disk. The fixed regular collar strips remain free of
equilibria. Proposition~\ref{rw:collar} places every counted
cycle essentially in the collar, surrounding that disk. The upper
vertex calculation in that proposition leaves at most one additional
equilibrium in the upper chart, and it is a saddle.
It remains to exclude horizontal equilibria from the cycle's disk.

If $\rho=0$, then $B=m=a=c=0$ and $\dot x=-y$; the
equilibrium equations give only the origin. Suppose $\rho>0$.
In a horizontal chart $x=\sigma/r$, $y=v/r$, $\sigma=\pm1$,
the equation $\dot y=0$ is
\[
 \sigma(r+v)+dv^2+crv=0.
\]
Its unique small solution is $v=rh(r)$ with $h=-1+o(1)$.
Thus $|z|=|1+y|<\delta$ at any horizontal equilibrium.
The other equilibrium equation reads
\[
 1-z+\bar BX^2+\bar aX+\epsilon\bar m(1-z)^2=0.
\]
Choose fixed small $R,\delta>0$ and then small $\epsilon$.
The left side is positive for $|X|\le R$, $|z|\le\delta$,
so the equilibrium lies in one of the two half-strips $|X|>R$.

In the incoming half-strip put $T=1/(-X)$ and use the positive
time change $\dd\tau/\dd\xi=-X$. With
$\eta_\ell=\rho\ell$, the exact equations are
\begin{equation}\label{rw:fibre:horizontal-equilibrium-field}
\begin{split}
 T_\tau&=\epsilon T\bigl[\bar B-\bar aT+
        \{1-z+\epsilon\bar m(1-z)^2\}T^2\bigr],\\
 z_\tau&=-z+\epsilon\bar cTz(z-1)
                             +\eta_\ell T(1-z)^2 .
\end{split}
\end{equation}
On $0\le T\le1/R$, $|z|\le\delta$, the tangential derivatives
are $O(\epsilon)$ and the normal derivative is
$-1+O(\epsilon+|\eta_\ell|)$.
Extend to a slightly larger rectangle, preserving $T=0$ and
these strict margins. The normally contracting graph construction
used in Section~4 gives
\[
 z=g(T)=\eta_\ell T\,O_{C^2}(1),\qquad
 Z=z-g(T),\qquad Z_\tau=-\beta(T,Z)Z,\quad\beta>1/2.
\]
The two factors in $g$ follow by solving on $\eta_\ell=0$
and $T=0$, respectively, and dividing these vanishing factors
in the bounded graph family. Every equilibrium lies on $Z=0$.
The small original quantity is $\rho\ell$; no bound on
$\nu=\ell/\rho$ is used.

A counted cycle has $z>0$ by
Proposition~\ref{rw:geometry:cycle-height}. At $T=1/R$,
$0\le z\le\delta$, one has $T_\tau>0$, while
$|g|<\delta/2$. A cycle point with $Z<0$ cannot leave the
rectangle in backward time through $T=1/R$; it cannot reach
$z=\delta$ because $Z<0$ is preserved and $g<\delta/2$.
It cannot reach $z=0$, and its compact affine orbit stays away
from $T=0$. It would therefore remain in the rectangle for
all backward time, whereas $|Z|$ grows exponentially.
This is impossible. A periodic orbit on $Z=0$ would be a
nonconstant periodic orbit of the scalar equation on that graph,
also impossible. Every cycle point in this half-strip consequently
lies strictly above $g$.

An equilibrium with $0<z=g(T)<\delta$ is joined vertically to
$z=0$ without meeting the cycle: cycle points with that $T$
are above $g$, including any points above the rectangle.
The line $z=0$ is outside the cycle's disk, since the whole
cycle has $z>0$. Thus the equilibrium is outside the disk.
The same holds immediately if its height is nonpositive.
Reflecting the field and reversing time gives the argument on
the outgoing half-strip, for either sign of $\ell$.
The coordinate $T$ includes finite roots and roots escaping
to infinity, so no horizontal escape case remains.

The single upper saddle is now the only possible additional
equilibrium inside the disk. Its index would change the enclosed total
from one to zero, contradicting the index of a periodic orbit. Thus
only the origin is enclosed. Two counted cycles are nested essential
curves and bound an equilibrium-free annulus in the collar. Cut this
annulus at the transverse central section. The compact-limit argument
for the strips below shows that every intermediate section point
returns, placing both cycles in one interval of the actual return
domain.
\end{proof}

\begin{proposition}[Saturation of the physical interval]
\label{rw:saturation}
Assume \eqref{rw:fibre:original-forcing}, fix
$p_0\in[p_\sharp/8,p_\sharp/4]$, and let $\gamma$ be a counted
cycle with $p_\gamma>p_0$. Every input
$p\in[p_0,p_\gamma]$ has an actual upper return $q(p)$
and both actual lower central intersections. They satisfy
\[
 q'(p)>0,\qquad L_+'(p)<0,\qquad L_-'(q)<0,
 \qquad
 L_+(p),\ L_-(q(p))\ge s_\gamma>0.
\]
All three arcs, and the segment joining their central endpoints,
lie in the fixed collar. This statement holds at nonclosing inputs
as well as on periodic orbits.
\end{proposition}
\begin{proof}
At $p_0$ the fixed small-endpoint upper map and both regular
lower maps exist. After reducing $p_\sharp$, their two actions
are at most $p_\sharp^2/31$ uniformly in the stated interval
for $p_0$. Their positive homogeneous amplitudes are bounded
below by a fixed positive prefactor times
$e^{-p_\sharp^2/(31\epsilon)}$.
The actual signed reference contributions are $O(|\ell|)$.
Since the original generator norm is bounded and the forcing
exponent is $p_\sharp^2/20$, both complete central heights are
positive for sufficiently small scale. Their arcs lie in the collar
by the regular construction.

The axis endpoints of the three initial arcs lie inside the
corresponding endpoints of $\gamma$. Since trajectories cannot cross
$\gamma$, the arcs stay in its disk. Join their central endpoints by
the central crosscut. Each lower path stays in its own half-plane until
its first central intersection, so it meets the crosscut only there.
The three arcs and the crosscut form a simple essential curve in the
collar, inside $\gamma$, whether or not its two central values coincide.

The initial and outer upper arcs, joined by the two axis segments,
bound a compact equilibrium-free strip by
Lemma~\ref{rw:fibre:enclosed-equilibria}. Since $\dot y=x$ on
$y=0$, the positive axis segment is an entrance, the negative one
an exit, and the orbit sides cannot be crossed. A nonexiting orbit
would have an equilibrium-free compact $\omega$-limit set, hence a
periodic orbit by Poincar\'e--Bendixson. Its Jordan disk lies in
$y>0$ inside $\gamma$ and contains no equilibrium, contradicting
the periodic-orbit index. Thus every intermediate upper entry returns, with
$q_0\le q(p)\le q_\gamma$. The transverse flow determinant
has a fixed sign and gives $q'(p)>0$.

Construct the right and left lower strips from the corresponding
initial and outer arcs, an axis segment, and a central crosscut.
They are equilibrium-free as well. On the right use backward
time, so the positive axis segment is an entrance; on the left
use forward time and the negative axis segment is an entrance.
The central crosscuts are exits because $\dot x>3/4$ on the
fixed central rectangle. The same compact-limit and index argument
gives both central intersections at every intermediate endpoint.
Transversality and endpoint order give $L_+'(p)<0$, $L_-'(q)<0$;
their minimum heights occur on $\gamma$ and equal $s_\gamma>0$.

The three strips and the joining central segments lie in the annulus
between the two essential curves, hence remain in the fixed collar.
The nonexit contradiction uses compactness for each fixed field;
flight times may diverge as a saddle or slow gap tends to zero.
\end{proof}

On this physical interval keep the actual central height as output.
The fixed weak references give
\[
 F(p)=L_+(p)-L_-(q(p))=U(p)-V(p)-\rho\nu H,\qquad H>0,
\]
where $H$ is independent of $p$ and $U',V'<0$. Consequently
\begin{equation}\label{rw:fibre:physical-log}
 K=\epsilon\log(V'/U'),\qquad
 F'=U'(1-e^{K/\epsilon}).
\end{equation}
A common regular input change and a change of the fixed additive
references preserve $K$ on this orbit interval. We keep the original
central height as output, since a nonlinear output change need not
preserve it. The signed center-comparison domains used to estimate
these functions are auxiliary; the physical interval of
Proposition~\ref{rw:saturation} has positive complete central heights.
